\section{Evaluasi}
\label{sec:evaluasi-bab8}

\subsection{Soal Pilihan Ganda}

\begin{enumerate}
    \item Mode pengalamatan yang dipakai oleh operand
          \lstasm|MOV AX, [BX+SI+4]| adalah \dots
          \begin{enumerate}
              \item \eng{indexed}
              \item \eng{base-relative}
              \item \eng{based-indexed}
              \item \eng{register indirect}
          \end{enumerate}

    \item Segmen \eng{default} untuk operand \lstasm|MOV AX, [BP+2]|
          adalah \dots
          \begin{enumerate}
              \item \reg{CS} \item \reg{DS}
              \item \reg{ES} \item \reg{SS}
          \end{enumerate}

    \item Penulisan operand berikut yang ditolak assembler karena
          tidak memiliki \eng{encoding} 8086 adalah \dots
          \begin{enumerate}
              \item \lstasm|MOV AX, [SI+DI]|
              \item \lstasm|MOV AX, [BX+DI]|
              \item \lstasm|MOV AX, [BP+DI]|
              \item \lstasm|MOV AX, [BX+SI]|
          \end{enumerate}

    \item Di antara pasangan instruksi berikut, yang paling cepat
          dieksekusi adalah \dots
          \begin{enumerate}
              \item \lstasm|MOV AX, [BX+SI+2]|
              \item \lstasm|MOV AX, [BP+2]|
              \item \lstasm|MOV AX, [1234h]|
              \item \lstasm|MOV AX, BX|
          \end{enumerate}

    \item Diketahui \reg{BX} = \hex{0100} dan \reg{SI} =
          \hex{0010}. Alamat efektif operand
          \lstasm|MOV AX, [BX+SI+8]| adalah \dots
          \begin{enumerate}
              \item \hex{0110} \item \hex{0118}
              \item \hex{0108} \item \hex{0128}
          \end{enumerate}
\end{enumerate}

\subsection{Soal Uraian}

\begin{enumerate}
    \item Jelaskan mengapa 8086 hanya mengizinkan \reg{BX},
          \reg{BP}, \reg{SI}, dan \reg{DI} sebagai penunjuk alamat
          serta mengapa \code{[AX]} dan \code{[SI+DI]} tidak dapat
          dikodekan. Kaitkan dengan byte \eng{ModR/M} pada
          Bab~\ref{chap:format-instruksi}.

    \item Tuliskan satu operasi yang sama, yaitu menyalin elemen
          berindeks pertama dari sebuah array \dir{DW} ke register
          \reg{AX}, dalam lima mode pengalamatan berbeda. Untuk
          setiap cara, sebutkan segmen \eng{default} yang dipakai.

    \item Diketahui \reg{DS} = \hex{1500}, \reg{SS} = \hex{2500},
          \reg{BX} = \hex{0040}, \reg{BP} = \hex{0060}, dan
          \reg{DI} = \hex{0006}. Hitung alamat efektif dan alamat
          fisik untuk \lstasm|MOV AX, [BX+DI+2]| dan
          \lstasm|MOV AX, [BP]|. Sebutkan pula langkah verifikasi yang
          Anda lakukan pada jendela memori \eng{debugger}.

    \item Tabel nilai mahasiswa berisi \eng{record} enam byte: empat
          byte NIM dan dua byte nilai. Rancang mode pengalamatan untuk
          membaca medan nilai mahasiswa ke-$i$, tuliskan potongan
          kodenya, dan jelaskan mengapa \lstasm|[BX+SI+4]| lebih
          tepat daripada tiga instruksi \lstasm|MOV| terpisah
          (Bab~\ref{chap:prosedur}).
\end{enumerate}

\subsection{Kisi-kisi Penilaian}

\begin{table}[H]
    \centering
    \caption{Kisi-kisi evaluasi Bab~\ref{chap:mode-pengalamatan}.}
    \label{tab:kisi-bab8}
    \small
    \begin{tabular}{@{}p{1.5cm}p{5.4cm}p{3.0cm}c@{}}
        \toprule
        \textbf{Sub-CPMK} & \textbf{Indikator yang diukur} &
        \textbf{Bentuk soal} & \textbf{Bobot} \\
        \midrule
        2.2 & Mengenali mode pengalamatan dari penulisan operand &
        Pilihan ganda 1; Uraian 2 & 25\% \\
        \addlinespace
        2.2 & Menghitung alamat efektif dan alamat fisik &
        Pilihan ganda 5; Uraian 3 & 25\% \\
        \addlinespace
        2.2 & Menentukan segmen \eng{default} dan
        \eng{override} & Pilihan ganda 2; Uraian 3 & 20\% \\
        \addlinespace
        2.2 & Menganalisis legalitas operand dan efisiensi mode &
        Pilihan ganda 3, 4; Uraian 1 & 15\% \\
        \addlinespace
        2.2 & Menerapkan mode pengalamatan pada array dan
        \eng{record} & Uraian 4 & 15\% \\
        \bottomrule
    \end{tabular}
\end{table}

\begin{checklistbox}
\textbf{Checklist Pencapaian Kompetensi Bab~\ref{chap:mode-pengalamatan}}

Seorang mahasiswa dinyatakan kompeten pada Sub-CPMK 2.2 apabila:
\begin{itemize}
    \item[\checked] dapat menyebutkan tujuh mode pengalamatan beserta
          contoh dan kegunaannya;
    \item[\checked] dapat menentukan segmen \eng{default} setiap mode
          serta menuliskan \eng{segment override prefix};
    \item[\checked] dapat menghitung alamat efektif dan alamat fisik
          operand memori tanpa kesalahan;
    \item[\checked] dapat menjelaskan alasan penolakan operand
          \code{[AX]}, \code{[BX+BP]}, dan \code{[SI+DI]};
    \item[\checked] dapat memilih mode yang efisien untuk akses array
          dan \eng{record};
    \item[\checked] dapat memverifikasi perhitungan alamatnya melalui
          jendela memori \eng{debugger}.
\end{itemize}
\end{checklistbox}
